\section*{ЗАКЛЮЧЕНИЕ}
\phantomsection
\addcontentsline{toc}{section}{Заключение}

В настоящей курсовой работе выполнен численный анализ влияния цилиндрического дискретного микрорельефа на рабочие характеристики радиального подшипника скольжения. В процессе исследования решены следующие задачи:

\begin{enumerate}
    \item Проведён обзор научной литературы по гидродинамической теории смазки и современным подходам к лазерному текстурированию поверхностей трения. Выявлено, что основным механизмом положительного воздействия микрорельефа является микрогидродинамический эффект, обусловленный формированием локальных клиновых зазоров на кромках углублений.

    \item Освоена методика численного решения уравнения Рейнольдса для подшипника конечной длины на базе метода конечных разностей с итерационной процедурой Гаусса--Зейделя и верхней релаксацией. Расчёт выполнен на сетке $500 \times 500$ узлов с критерием сходимости $\delta < 10^{-5}$.

    \item Определены поля давления и интегральные характеристики (нагрузочная способность $F$, коэффициент трения $\mu$, расход смазки $Q$) для гладкого подшипника и подшипника с цилиндрическими углублениями (тип~3) в диапазоне эксцентриситетов $\varepsilon = 0{,}05 - 0{,}80$. Геометрические параметры подшипника: $R = 50$~мм, $c = 70$~мкм, $L = 80$~мм (геометрия~B).

    \item Выполнено сравнение характеристик гладкого и текстурированного подшипников. При характерном рабочем эксцентриситете $\varepsilon = 0{,}6$ зафиксированы следующие изменения:
    \begin{itemize}
        \item нагрузочная способность увеличилась на $13{,}05$\%: с $F = 10685{,}6$~Н до $F = 12080{,}2$~Н;
        \item коэффициент трения снизился на $12{,}47$\%: с $\mu = 0{,}01507$ до $\mu = 0{,}01319$;
        \item расход смазки возрос на $0{,}72$\%: с $Q = 0{,}6739$~л/с до $Q = 0{,}6788$~л/с.
    \end{itemize}
\end{enumerate}

Полученные данные свидетельствуют о том, что нанесение цилиндрических углублений с параметрами $a = 2{,}4$~мм, $b = 2{,}2$~мм, $h_p = 10$~мкм ($h_p/c = 0{,}1$) приводит к заметному улучшению трибологических характеристик подшипника. Совместное повышение нагрузочной способности и уменьшение коэффициента трения при минимальном росте расхода смазки свидетельствует о практической целесообразности данного типа текстурирования.

Цилиндрический профиль углублений отличается постоянной глубиной внутри границ и резкими переходами на кромках, что формирует интенсивные локальные градиенты давления в смазочном слое. Подшипник геометрии~B с увеличенными габаритами ($D = 100$~мм, $L = 80$~мм) обеспечивает значительно более высокие абсолютные значения нагрузочной способности по сравнению с геометрией~A, хотя относительный эффект текстурирования при данных параметрах углублений выражен умереннее. Тем не менее прирост несущей способности свыше 13\% и снижение трения на 12{,}5\% представляют существенное улучшение эксплуатационных характеристик.

На основании проведённого исследования можно рекомендовать применение цилиндрического микрорельефа для подшипников скольжения нефтегазового оборудования, работающего в условиях жидкостного трения. Практический интерес представляет использование текстурирования в тяжелонагруженных узлах ($\varepsilon > 0{,}5$), где абсолютный прирост нагрузочной способности наиболее значителен, а снижение потерь на трение способствует повышению энергоэффективности оборудования.
